% =============================================================================
% KAPITEL 2: GRUNDLAGEN
% =============================================================================

\chapter{Grundlagen}

In diesem Kapitel werden die Grundlagen erklärt.

% -----------------------------------------------------------------------------
% Erster Abschnitt
% -----------------------------------------------------------------------------
\section{Erster Abschnitt}

\blindtext

\subsection{Unterabschnitt}

\blindtext

\begin{figure}[htbp]
	\centering
	%\includegraphics[width=0.8\textwidth]{beispielbild}
	\caption[Beispielabbildung]{Dies ist eine Beispielabbildung}
	\label{fig:beispiel}
\end{figure}

Wie in \autoref{fig:beispiel} zu sehen ist, handelt es sich um ein wichtiges Konzept.

% -----------------------------------------------------------------------------
% Zweiter Abschnitt
% -----------------------------------------------------------------------------
\section{Zweiter Abschnitt}

\blindtext

% Beispiel für eine Tabelle
\begin{table}[htbp]
	\centering
	\caption{Beispieltabelle}
	\begin{tabular}{ll}
		\toprule
		\textbf{Spalte 1} & \textbf{Spalte 2} \\
		\midrule
		Wert 1            & Wert 2            \\
		Wert 3            & Wert 4            \\
		\bottomrule
	\end{tabular}
	\label{tab:beispiel}
\end{table}

\autoref{tab:beispiel} zeigt ein Beispiel für eine Tabelle.

% -----------------------------------------------------------------------------
% Code-Beispiel
% -----------------------------------------------------------------------------
\section{Code-Beispiel}

Hier ein Beispiel in Python:

\lstset{style=Python}
\begin{lstlisting}
def hallo_welt():
    print("Hallo, Welt!")
    
hallo_welt()
\end{lstlisting}

% Zurück zum Standard-Stil
\lstset{language=[LaTeX]TeX}
